\documentclass[runningheads,11pt]{llncs}

\usepackage[T1]{fontenc}
\usepackage{lmodern}
\usepackage{amsmath,amssymb}
\usepackage{alltt}
\usepackage{booktabs}
\usepackage[hidelinks]{hyperref}

\DeclareUrlCommand\lean{\urlstyle{tt}}
\newcommand{\efxz}{\texorpdfstring{\mbox{EFX$_0$}}{EFX0}}
\newcommand{\LBp}{\texorpdfstring{\mbox{LB$^{+}$}}{LB+}}
\newcommand{\NA}{\mathit{NA}}
\newcommand{\capa}{\mathrm{cap}}
\newcommand{\base}{\mathrm{base}}
\newcommand{\none}{\bot}
\newcommand{\st}{\mathrel{:}}

\begin{document}

\title{Complete \efxz{} Allocations Exist When Every Agent Values at Most Three Goods}
\titlerunning{\efxz{} with at Most Three Relevant Goods}

\author{Evan Lin}
\authorrunning{E. Lin}
\institute{University of Illinois Urbana-Champaign}

\maketitle

\begin{abstract}
An allocation of indivisible goods is \efxz{} if no agent envies another agent's bundle after the removal of any single good from it, even a good that the envious agent values at zero. We prove that every instance with nonnegative additive valuations in which each agent positively values at most three goods has a complete \efxz{} allocation, for any number of agents. The proof is constructive. After agents that value at most two remaining goods or have a dominant good are peeled off, a serial dictatorship with a careful placement of the unpicked goods, followed if needed by one rotation of picks along a chain of agents, gives an \efxz{} allocation in which at most one bundle has more than two goods. The resulting algorithm uses at most $400(n+m+1)^4$ counted operations. The existence theorem, the structural theorem, the algorithm and its running-time bound are machine-checked in Lean~4, with core Lean and the standard axioms only. Bundles of at most two goods do not always suffice.
\keywords{Fair division \and Indivisible goods \and EFX \and Serial dictatorship \and Formal verification \and Lean}
\end{abstract}

\section{Introduction}

Envy-freeness cannot always be achieved with indivisible goods, and envy-freeness up to any good is one of its most studied relaxations~\cite{AmanatidisBFV22}. Let agents $N=[n]$ share goods $M=[m]$, with additive valuations $v_i(S)=\sum_{g\in S}v_i(g)$, $v_i(g)\ge0$. An allocation $X=(X_1,\dots,X_n)$, a partition of $M$, is \emph{\efxz{}} if
\[
  v_i(X_i)\ \ge\ v_i(X_j\setminus\{h\})\qquad\text{for all } i\ne j \text{ and all } h\in X_j,
\]
and \emph{EFX} if this is required only for goods $h$ with $v_i(h)>0$; the stronger notion is usually called \efxz{}~\cite{AmanatidisBFV22}. The difference matters when agents value few goods: if $X_j$ contains a good worthless to $i$, then \efxz{} requires $v_i(X_i)\ge v_i(X_j)$. In particular, goods that nobody values cannot be handed to an arbitrary agent.

Complete EFX allocations are known to exist for additive valuations only in special cases, among them three agents~\cite{ChaudhuryGM20}, $m\le n+3$~\cite{Mahara21}, four agents with at most nine goods~\cite{AlkassarFM26}, when every good is valued by at most two agents~\cite{ChristodoulouFKS23,AfshinmehrAMMS26,KavianiSS24,SgouritsaS25}, and on hypergraphs of girth at least four~\cite{LianeasSS26}; these are \efxz{} results. For monotone valuations EFX can fail~\cite{AkramiMMSW26,MackenzieS26}. Call an instance \emph{$k$-relevant} if every agent positively values at most $k$ goods. For $k\le2$, serial dictatorship gives an \efxz{} allocation~\cite{MrdEfx}. For $k\le4$, Viswanathan and Mehta~\cite{ViswanathanM24} announce ordinary EFX, with a proof sketch; their algorithm gives the goods nobody values to an arbitrary agent, which can break \efxz{}. We ask whether every $3$-relevant instance has a complete \efxz{} allocation. A counterexample would refute the additive EFX conjecture, since perturbing all values by a small $\varepsilon>0$ (compare~\cite[Lemma~1]{ChaudhuryGM20}) turns an instance without \efxz{} allocations into a positive one without EFX allocations. But the perturbation destroys the relevance structure, and we know of no result that implies a positive answer (the literature search is recorded in~\cite{Repo}). We prove:

\begin{theorem}[existence]\label{thm:target}
Every instance with $n\ge1$ agents and nonnegative real additive valuations in which every agent positively values at most three goods has a complete \efxz{} allocation.
\end{theorem}

\noindent Call an agent \emph{balanced} if it values exactly three goods and none of them more than the other two together.

\begin{theorem}[one large bundle]\label{thm:D}
In every instance with $n\ge1$ agents and nonnegative real additive valuations in which every agent is balanced, some \efxz{} allocation has at most one bundle with more than two goods.
\end{theorem}

\noindent The construction, \LBp{} (Section~\ref{sec:lbplus}), is a serial dictatorship with a careful placement of the unpicked goods and at most one \emph{rotation} of picks along a chain of agents. It gives a polynomial-time algorithm (Section~\ref{sec:algo}), and everything is machine-checked in Lean~4 (Section~\ref{sec:lean}). A large bundle is sometimes necessary, and may need four goods (Proposition~\ref{prop:neg}).

\paragraph{Status of the claims.} The work was done in a public repository~\cite{Repo} whose ledger gives every claim a status: PROVED (complete written proof), CERTIFIED (exhaustive computation over a provably complete finite set, with independently checked certificates), EVIDENCE, CONJECTURE or REFUTED. All positive results below are PROVED. Two independent reviews of the written proof of Theorems~\ref{thm:target} and~\ref{thm:D} found no error.

\section{Preliminaries}\label{sec:prelim}

Let $R_i=\{g\st v_i(g)>0\}$. Since $\max_{h\in X_j}v_i(X_j\setminus\{h\})=v_i(X_j)-\min_{h\in X_j}v_i(h)$, a bundle of one good is never strongly envied, and $i$ must not envy at all a bundle that contains a good outside $R_i$. If $|R_i|=3$ we write $R_i=\{a_i,b_i,c_i\}$ with $v_i(a_i)\ge v_i(b_i)\ge v_i(c_i)>0$, fix this order once (ties broken arbitrarily) as $a_i\succ_i b_i\succ_i c_i$, and call $i$ \emph{balanced} if $v_i(a_i)\le v_i(b_i)+v_i(c_i)$, \emph{strictly balanced} if the inequality is strict, and \emph{top-heavy} if $v_i(a_i)\ge v_i(b_i)+v_i(c_i)$.

\begin{lemma}[peeling]\label{lem:peel}
Let $i\in N$ and $P\subseteq M$ satisfy (i)~$v_i(P)\ge v_i(M\setminus P)$ and (ii)~$|P|\le1$ or $P$ is worthless to every other agent. If the instance without $i$ and $P$ has an \efxz{} allocation $Y$, then $Y$ together with $X_i=P$ is \efxz{}.
\end{lemma}

\noindent \emph{Rule R1} applies Lemma~\ref{lem:peel} with $P=\{p\}$, $p$ a favourite remaining good of $i$, when $v_i(p)\ge v_i(M\setminus\{p\})$ (that is, $i$ values at most two remaining goods or is top-heavy), and with $P=\emptyset$ when $i$ values no remaining good. \emph{Rule R2} gives $i$ its goods that no other agent values, if there are at least two and (i) holds. Goods valued by nobody go to an unenvied agent after envy cycles are rotated away (Lemma~\ref{lem:junk}). What remains when none of these applies is a single agent or a \emph{core}: at least two agents, each valuing exactly three goods, strictly balanced, with at most one private good, and no good valued by nobody.

\begin{proposition}[core reduction]\label{prop:core}
If every core with at most $N_0$ agents has an \efxz{} allocation, so does every $3$-relevant instance with between $1$ and $N_0$ agents.
\end{proposition}

\noindent Cores are the objects of our computations (Section~\ref{sec:lean}). For Theorem~\ref{thm:target}, rule R1 alone suffices: if R1 fails for every agent's favourite good $p$, then $2v_i(g)\le2v_i(p)<v_i(M)$ for all $g$, so every agent of a $3$-relevant instance values exactly three goods and is strictly balanced, and Theorem~\ref{thm:D} allows goods valued by nobody and several private goods.

\begin{lemma}[real values reduce to naturals]\label{lem:L12}
If $v$ is nonnegative real additive with $|R_i|\le3$ for all $i$, there are natural-number additive valuations $w$ with the same relevant goods such that $v_i(S)\le v_i(T)\iff w_i(S)\le w_i(T)$ for every agent $i$ and all $S,T\subseteq M$.
\end{lemma}

\noindent Such comparisons are fixed by the ties among $a_i,b_i,c_i$ and the sign of $v_i(a_i)-v_i(b_i)-v_i(c_i)$, and \efxz{}, relevance and balance are made of them. Proofs are in Appendix~\ref{app:prelim}.

\section{The Construction \LBp{}}\label{sec:lbplus}

In this section every agent values exactly three goods and is balanced; goods valued by nobody and agents with several private goods are allowed.

\paragraph{Phase 1: serial dictatorship with R1 priority.} Agents are processed one at a time; $G$ is the set of goods not yet picked, initially $M$. The processed agent $i$ takes its $\succ_i$-favourite good of $R_i\cap G$ as its \emph{pick} $Y_i$, or no pick ($Y_i=\none$) if $R_i\cap G=\emptyset$. If some unprocessed agent has $|R_i\cap G|\le2$, the next agent is any such agent (an \emph{R1 step}). Otherwise the next agent is any unprocessed agent, and it takes its top (an \emph{insertion step}). The goods never picked form the \emph{junk} $J_0$. A \emph{block} is an insertion step with the R1 steps that follow it; its inserted agent is its \emph{leader}. We use four immediate invariants (Appendix~\ref{app:proofs}): (I1)~if $g\in R_i$ is above $Y_i$ (or $Y_i=\none$), then $g$ was picked before $i$'s turn, and (B)~by an agent of $i$'s block; (I2)~if $g\in R_i\cap J_0$, then $Y_i\succ_i g$; (I3)~an agent that finds all its goods available is a leader.

\paragraph{Pre-allocations.} A \emph{pre-allocation} $P=(Y,U)$ has picks $Y_i\in R_i\cup\{\none\}$ (distinct) and a set $U$ of \emph{upgraded} agents; each $u\in U$ has $Y_u=b_u$ and also holds $c_u$, and these goods $c_u$ are nobody's pick. The \emph{junk} $J$ is the set of the other goods. For $i\notin U$ the \emph{needs} are $N_i=\{g\in R_i\st g\succ_i Y_i\}$ ($N_i=R_i$ if $Y_i=\none$), and $\NA=\bigcup_{i\notin U}N_i$ is the set of goods \emph{needed alone}. $P$ is \emph{valid} if (V1)~$J\cap\NA=\emptyset$ and (V2)~$b_u,c_u\notin\NA$ for $u\in U$. An agent $i\notin U$ is \emph{frozen} if $Y_i\in\NA$ and a \emph{terminal} otherwise. The \emph{slots} are $\capa(i)=0$ for upgraded and frozen agents, and $1$ or $2$ for a terminal with or without a pick; $S=\sum_i\capa(i)$ and $\omega=|J|-S$. The \emph{base} of $i$ is $\{Y_i\}$ (or $\emptyset$) for $i\notin U$ and $\{b_i,c_i\}$ for $i\in U$. A \emph{completion} with \emph{owner} $o$ (a terminal or upgraded agent) gives every agent its base, at most $\capa(i)$ junk goods to each terminal $i\ne o$, and the remaining junk to $o$; without owner, all junk goes into slots. It satisfies the \emph{owner constraint} (OC) if no agent $x\notin U$, $x\ne o$, with $Y_x=a_x$ has $\{b_x,c_x\}\subseteq X_o$.

\begin{theorem}[soundness]\label{thm:sound}
Let $P$ be valid. Every completion of $P$ that satisfies (OC) is \efxz{} for every balanced additive valuation consistent with the rankings, and only the owner's bundle can have more than two goods.
\end{theorem}

\begin{proof}
Only $X_o$ can exceed two goods. Fix $i$ and $j\ne i$ with $|X_j|\ge2$. Then $j$ is not frozen, so $X_j\cap\NA=\emptyset$ by the definition of frozen agents, (V1) and (V2). If $i\in U$, then $X_j$ meets $R_i$ in at most $a_i$, and $v_i(X_j\setminus\{h\})\le v_i(a_i)\le v_i(b_i)+v_i(c_i)\le v_i(X_i)$. If $Y_i=\none$, then $R_i=N_i\subseteq\NA$ and $X_j$ is worthless to $i$. Otherwise every good of $R_i\cap X_j$ is outside $N_i\cup\{Y_i\}$, hence below $Y_i$. If there is at most one, $v_i(X_j\setminus\{h\})\le v_i(Y_i)\le v_i(X_i)$. If there are two, then $Y_i=a_i$ and $R_i\cap X_j=\{b_i,c_i\}$; for $|X_j|=2$ removing one leaves a good worth at most $v_i(a_i)$, and $|X_j|\ge3$ means $j=o$, excluded by (OC).\qed
\end{proof}

\noindent The \emph{exposed} agents of an owner $o$ are $E_o=\{x\notin U,\ x\ne o\st Y_x=a_x,\ \{b_x,c_x\}\subseteq J\cup\base(o)\}$. By filling the other slots with the goods of $H$ first, one gets:

\begin{lemma}[owner criterion]\label{lem:owner}
Let $P$ be valid, $\omega\ge1$. If no $x\in E_o$ has $\{b_x,c_x\}\subseteq\base(o)$, and some $H\subseteq J$ with $|H|\le S-\capa(o)$ meets $\{b_x,c_x\}$ for every $x\in E_o$, then $P$ has a completion with owner $o$ satisfying (OC).
\end{lemma}

\noindent Then $o$ is a \emph{valid owner}. If $\omega\le0$, the completion without owner has all bundles of at most two goods.

\paragraph{Upgrades and the natural owner.} After Phase 1, let $U=\emptyset$; while some $k\notin U$ has $Y_k=b_k$, $c_k\in J$ and $b_k\notin\NA$, add $k$ to $U$ (in any order). The result is valid, by (I2) and because $\NA$ only shrinks (Lemma~\ref{lem:upgrades}). The construction LB of~\cite{Repo} now searches for any valid owner and fails if there is none. \LBp{} tries the last-processed agent $r\notin U$, a terminal. Every frozen agent $x$ has a \emph{need chain} $x=x_0,\dots,x_t$ of distinct agents in which each $x_{s+1}\notin U$ needs $Y_{x_s}$, $x_0,\dots,x_{t-1}$ are frozen and $x_t$ is a terminal; by (B) it moves forward inside $x$'s block.

\begin{theorem}[the owner $r$]\label{thm:A}
Let $\omega\ge1$ after Phase 1 and the upgrades. Then $r$ is a valid owner unless the \emph{bad case} holds: an agent $k^*$ of $r$'s block is exposed for $r$ and frozen, every need chain from $k^*$ ends at $r$, and the sets $\pi_x=\{b_x,c_x\}\cap J$ ($x\in E_r$) are pairwise disjoint.
\end{theorem}

\begin{proof}[idea]
An exposed agent found all its goods available ($b_x,c_x$ are junk or $r$'s later pick), so by (I3) it leads its block, and each block has at most one. For each exposed $x$ outside $r$'s block, $x$ or the end of a need chain from $x$ is a terminal of $x$'s block. These terminals are distinct and differ from $r$, so $S-\capa(r)\ge|E_r|-1$, and $\ge|E_r|$ if no $k^*$ exists. One junk good per set $\pi_x$, sharing a good where two sets meet, fits unless the bad case holds (Appendix~\ref{app:proofs}).\qed
\end{proof}

\begin{theorem}[rotation]\label{thm:B}
In the bad case, take any need chain $k^*=x_0,\dots,x_t=r$. Let each $x_s$ ($s\ge1$) take $Y_{x_{s-1}}$, and let $k^*$ take $b_{k^*},c_{k^*}$ and join $U$ ($Y_r$ is released). The result $P'$ is valid, and if $\omega'\ge1$, then $k^*$ is a valid owner of $P'$.
\end{theorem}

\begin{proof}[idea]
Chain agents receive goods they needed, so $\NA$ shrinks; $b_{k^*},c_{k^*}\in J\cup\{Y_r\}$ are free. Every agent exposed for $k^*$ in $P'$ was exposed for $r$ (not the interior of the chain, and not $r$, by (I1) and the stopping rule of the upgrades), and disjointness keeps its pair out of $\{b_{k^*},c_{k^*}\}$. The terminals found for the other exposed agents lie outside $r$'s block and keep their slots, which suffice (Appendix~\ref{app:proofs}).\qed
\end{proof}

\paragraph{Construction \LBp{}.} Run Phase~1 (any choices) and the upgrades (any order). If $\omega\le0$, return the completion without owner. Else, if $r$ is a valid owner, return a completion with owner $r$. Else, by Theorem~\ref{thm:A}, the bad case holds: rotate along any need chain from $k^*$ to $r$ and return a completion of $P'$, without owner if $\omega'\le0$ and with owner $k^*$ otherwise.

\begin{theorem}[\LBp{} never fails]\label{thm:C}
If every agent values exactly three goods and is balanced, then for every choice in Phase~1, in the upgrades and of the chain, \LBp{} returns an allocation that is \efxz{} for every balanced additive valuation consistent with the rankings, with at most one bundle of more than two goods.
\end{theorem}

\begin{proof}
The pre-allocation is valid, and the completion satisfies (OC), by Lemma~\ref{lem:owner} and Theorems~\ref{thm:A} and~\ref{thm:B}. Apply Theorem~\ref{thm:sound}.\qed
\end{proof}

\begin{proof}[of Theorems~\ref{thm:target} and~\ref{thm:D}]
Theorem~\ref{thm:D} is Theorem~\ref{thm:C}. For Theorem~\ref{thm:target}, induct on $n$: one agent takes everything; if R1 applies to an agent, peel it (Lemma~\ref{lem:peel}); otherwise every agent is balanced (Section~\ref{sec:prelim}) and Theorem~\ref{thm:D} applies.\qed
\end{proof}

\noindent The proof uses no computation. The large bundle has $|\NA|-(2n-m)+2$ goods (a written remark, not formalized).

\begin{example}
Let agents $0,1,2$ rank $g_0\succ g_2\succ g_3$, $g_2\succ g_1\succ g_4$, $g_0\succ g_2\succ g_1$, processed in the order $1,0,2$. Agent~1 is inserted and picks $g_2$; agents $0$ and $2$ then pick $g_0$ and $g_1$ in R1 steps. Now $J=\{g_3,g_4\}$, $\NA=\{g_0,g_2\}$, agents $0,1$ are frozen, and $r=2$ has one slot: $\omega=1$. The bundle $\{g_1,g_3,g_4\}$ for $r$ violates (OC) for agent~1 ($b_1=g_1$, $c_1=g_4$), so LB fails on this run. The rotation along the chain $1\to2$ gives $\{g_0\},\{g_1,g_4\},\{g_2,g_3\}$, which is \efxz{}.
\end{example}

\section{Algorithm and Running Time}\label{sec:algo}

\emph{Algorithm K3ALG.} Stage~R: while at least two agents and some goods remain, remove the first agent to which rule R1 applies, with its favourite good or nothing; a lone remaining agent takes all goods left. Stage~L: if agents remain, each values exactly three remaining goods and is strictly balanced; sort each agent's goods and run \LBp{} with the processing order ``the first agent with at most two goods left, else the first agent'', upgrades by index, the owner test below, and a need chain found by scanning forward from $k^*$ (pseudocode in Appendix~\ref{app:algo}).

\begin{theorem}[algorithm]\label{thm:algo}
On every instance with $n\ge1$ agents and natural-number values in which every agent values at most three goods, K3ALG returns an \efxz{} allocation. On every instance with $n\ge1$ agents it performs at most $400(n+m+1)^4$ counted operations.
\end{theorem}

\noindent Correctness follows from Lemma~\ref{lem:peel} and Theorem~\ref{thm:C}. The one step that is not obviously polynomial is the owner test: a minimum set of junk goods meeting all sets $\pi_x$ (each of at most two goods) is a vertex cover. For $r$ it is not needed.

\begin{proposition}[the owner test]\label{prop:O}
With $\omega\ge1$ after Phase~1 and the upgrades, $r$ is a valid owner if and only if $|E_r|\le S-\capa(r)$ or two of the sets $\pi_x$ ($x\in E_r$) meet.
\end{proposition}

\begin{proof}
The proof of Theorem~\ref{thm:A} gives $S-\capa(r)\ge|E_r|-1$, and nonempty sets $\pi_x$ have a hitting set of fewer than $|E_r|$ goods if and only if two of them meet.\qed
\end{proof}

\noindent In Lean, K3ALG is a program in a \emph{cost monad} (a value paired with an operation count), and \lean{EFX.K3.algo} is defined as its value, so the bound is about the program that computes the allocation. One unit is a list step, a comparison of agents or goods, a read of an input value or of a table, or a comparison or addition of natural numbers of any size. The dominant term is the upgrade loop. Finer bounds, $O(n^4+n^2m)$ and $O((n+m)\log n)$ after reading the input for a Python implementation, are written but not formalized; that implementation agrees with the Lean program on every instance tested and takes about 1.6\,s with $10^5$ agents (evidence). K3ALG compares only subset sums of one agent's values, so by Lemma~\ref{lem:L12} it runs verbatim on real values (a written remark, not in the ledger).

\section{Machine-Checked Proof and Computational Certification}\label{sec:lean}

\paragraph{Lean.} The formalization uses core Lean~4 (v4.34.0) and no packages. No file contains \texttt{sorry}, the build is warning-free, and every declaration depends only on \texttt{propext}, \texttt{Classical.choice} and \texttt{Quot.sound}; a script checks this with a \verb|#print axioms| certificate per ledger-facing theorem, a whole-library axiom check and Lean's replay checker. The trusted base is six short definitions copied from~\cite{MrdEfx} (Appendix~\ref{app:lean}): values $v\colon\mathrm{Fin}\,n\to\mathrm{Fin}\,m\to\mathbb{N}$, an allocation as a map from goods to agents, and \efxz{} with the removed good ranging over \emph{every} good of the other bundle. Theorem~\ref{thm:target} is \lean{EFX.target}, Theorem~\ref{thm:D} is \lean{EFX.LB.corollaryD} (balance as $2v_i(g)\le v_i(M)$), and Theorem~\ref{thm:algo} is \lean{EFX.K3.algo_efx0} with \lean{EFX.K3.algoC_cost}. Via Lemma~\ref{lem:L12} (\lean{EFX.l12}), \lean{EFX.target_ordered} and \lean{EFX.corollaryD_ordered} extend the first two to values in any linearly ordered cancellative commutative monoid; that $\mathbb{R}_{\ge0}$ is one is the textbook fact, as core Lean has no reals. Every result of Section~\ref{sec:lbplus} is formalized, Theorem~\ref{thm:C} over all choices of \LBp{} (Appendix~\ref{app:lean}); the final statements use only the trusted base. An audit restated Theorems~\ref{thm:target} and~\ref{thm:D} independently, before reading the model, and derived them from the main theorems; a second kernel (nanoda) re-checked their natural-number development. A second Lean proof of Theorem~\ref{thm:D} for cores, through Pareto-maximal pre-allocations and independent of \LBp{}, is also in~\cite{Repo}. Not formalized: the size of the large bundle, the finer running-time bounds and the real-value remark on K3ALG.

\paragraph{Computation.} Small cases are also certified independently of the proof (Appendix~\ref{app:cert}). In a core, \efxz{} depends only on the agents' rankings (Lemma~\ref{lem:L5}), and for every connected core (enumerated by nauty's \texttt{genbg} and cross-checked) stored allocations cover all $6^n$ ranking profiles, as a SAT-free checker verifies from the raw \efxz{} definition. This certifies \efxz{} for all $3$-relevant instances with $n\le7$, and Theorem~\ref{thm:D} for all cores with $n\le6$ and all connected cores with $n=7$, or $n=8$ and $m\ge13$. Computation also refuted stronger shapes; Proposition~\ref{prop:neg} gives hand-checked counterexamples.

\section{Conclusion and Open Problems}\label{sec:concl}

Every $3$-relevant instance has a complete \efxz{} allocation, found in polynomial time, with a machine-checked proof. The case of four relevant goods per agent is open; the repository has work in progress on it, about which we claim nothing. Also open: whether LB, without rotation, never fails; the minimum size and canonical content of the large bundle (conjecture~K of~\cite{Repo}); and the complexity of the owner test for arbitrary owners.

\paragraph{Acknowledgments.} The proofs, code and Lean formalization were developed with the assistance of AI coding agents (Claude Code). Every claim in this paper is machine-checked, computer-certified or refereed, as stated at the claim.

\bibliographystyle{splncs04}
\bibliography{refs}

\clearpage
\appendix

\section{Proofs for Section~\ref{sec:prelim}}\label{app:prelim}

Full proofs of all reduction lemmas are in \texttt{proofs/lemmas.md} of~\cite{Repo}; we reproduce those used here.

\begin{proof}[of Lemma~\ref{lem:peel}]
Every other bundle lies in $M\setminus P$, so $i$ values it at most $v_i(M\setminus P)\le v_i(P)=v_i(X_i)$. Towards $P$, every other agent sees a single good or a bundle worthless to it. The remaining constraints are those of $Y$.\qed
\end{proof}

\begin{lemma}[junk]\label{lem:junk}
Let $J\ne\emptyset$ be the goods valued by no agent. If the instance without $J$ has an \efxz{} allocation, so does the instance.
\end{lemma}

\begin{proof}
Let $Y$ be \efxz{} for $M\setminus J$. While the envy graph ($i\to j$ iff $v_i(Y_j)>v_i(Y_i)$) has a cycle, give each agent on it the bundle of its successor. The bundles are the same and no agent loses value; an agent that now sees its old bundle satisfies $v_i(Y_i\setminus\{h\})\le v_i(Y_i)$, so \efxz{} is preserved. The sum of values strictly increases, so this stops, with an acyclic envy graph. Let $s$ be a source: $v_i(Y_s)\le v_i(Y_i)$ for all $i$. Put $X_s=Y_s\cup J$. For $i\ne s$, $v_i(X_s\setminus\{h\})\le v_i(X_s)=v_i(Y_s)\le v_i(Y_i)$; all other constraints are unchanged.\qed
\end{proof}

\begin{proof}[of Proposition~\ref{prop:core}]
Induction on $n+m$ over $3$-relevant instances with $1\le n\le N_0$. One agent takes everything. Otherwise: (1)~if some good is valued by nobody, use Lemma~\ref{lem:junk}; (2)~else if some agent has at most two relevant goods or is top-heavy, use rule R1 and Lemma~\ref{lem:peel}; (3)~else if some agent $i$ has at least two private goods, let $P$ be all of them: if $|P|=3$, (i) is trivial, and if $P=\{x,y\}$ and $R_i\setminus P=\{z\}$, then $v_i(x)+v_i(y)\ge v_i(b_i)+v_i(c_i)>v_i(a_i)\ge v_i(z)$, so Lemma~\ref{lem:peel} applies (rule R2); (4)~else the instance is a core, covered by hypothesis. Each case passes to a smaller $3$-relevant instance.\qed
\end{proof}

\noindent In Lean this is \lean{EFX.core_reduction} (with \lean{EFX.peel}, \lean{EFX.peelEmpty}, \lean{EFX.peelR2} and \lean{EFX.Inst.exists_efx0_of_junk}).

\begin{proof}[of Lemma~\ref{lem:L12}]
Fix $i$; $w_i$ is built from $v_i$ alone. Goods outside $R_i$ contribute $0$ to both sides under both valuations, and removing goods common to $S$ and $T$ changes neither comparison, so it suffices that $w_i$ preserves $v_i(X)\le v_i(Y)$ for disjoint $X,Y\subseteq R_i$. Name the relevant goods with values $a\ge b\ge c>0$. Comparisons with an empty side are decided by positivity; two single goods by the ties among $a,b,c$; a single good against the other two by the sign of $a-(b+c)$, because $b<a+c$ and $c<a+b$ always hold. Two disjoint pairs do not fit in three goods. Choose $w_i$ on $R_i$ (in the same order) by the pattern: $(4,3,2)$ if $a>b>c$ and $a<b+c$; $(3,2,1)$ if $a>b>c$ and $a=b+c$; $(5,2,1)$ if $a>b>c$ and $a>b+c$; $(2,2,1)$ if $a=b>c$; $(3,2,2)$, $(2,1,1)$, $(3,1,1)$ if $a>b=c$ and $a<2b$, $a=2b$, $a>2b$; $(1,1,1)$ if $a=b=c$; $(2,1)$ or $(1,1)$ for two relevant goods; $(1)$ for one; and $w_i(g)=0$ for $g\notin R_i$. Each triple has the same ties and the same sign of $a-(b+c)$ as its pattern.\qed
\end{proof}

\noindent Lean: \lean{EFX.l12}, which uses all 31 permutations of these representatives (checked by \texttt{decide}) instead of sorting, and \lean{EFX.efx0_iff_of_agree} for the transfer of \efxz{}.

\section{Proofs for Section~\ref{sec:lbplus}}\label{app:proofs}

These are the proofs of \texttt{proofs/lb\_last\_step.md} of~\cite{Repo}, Sections 1--6.

\paragraph{Phase 1.} (I1): $i$ took its favourite remaining good, and goods leave $G$ only by being picked. (I2): $g$ was still in $G$ at $i$'s turn. (I3): R1 steps process only agents with at most two goods left. (B): each block starts with an insertion step, when every unprocessed agent $i$ has $R_i\subseteq G$; by (I1) $g$ was picked before $j$'s turn, and not in an earlier block, since $j$ was unprocessed when its own block started.

\begin{proof}[of Lemma~\ref{lem:owner}]
Since $|J|=S+\omega>S-\capa(o)$, $H$ extends to $C\subseteq J$ with $|C|=S-\capa(o)$, which exactly fills the slots of the terminals other than $o$; let $X_o=\base(o)\cup(J\setminus C)$. If $x\notin U$, $x\ne o$, $Y_x=a_x$ had $\{b_x,c_x\}\subseteq X_o$, then $\{b_x,c_x\}\subseteq J\cup\base(o)$, so $x\in E_o$, and $H\subseteq C$ meets $\{b_x,c_x\}$ outside $X_o$, a contradiction.\qed
\end{proof}

\begin{lemma}[upgrades]\label{lem:upgrades}
When the upgrade loop stops, $(Y,U)$ is a valid pre-allocation, and (UT) no $k\notin U$ has $Y_k=b_k$, $c_k\in J$ and $b_k\notin\NA$.
\end{lemma}

\begin{proof}
(UT) is the stopping condition. (V1): by (I2) a junk good lies below the pick of every agent that values it, so it is in no $N_i$. (V2): $c_u$ was junk; $b_u\notin\NA$ held when $u$ was upgraded, and $\NA$ only shrinks as $U$ grows.\qed
\end{proof}

\begin{proof}[of Theorem~\ref{thm:A}]
$r$ exists: the first processed agent is inserted and holds its top, so it is not upgraded.
(A1)~$r\in T$: if $Y_r=\none$, $r$ is not frozen by definition; if $Y_r\in N_j$ for some $j\notin U$, then by (I1) $j$ was processed after $r$, contradicting the choice of $r$.
(A2)~Every agent processed after $r$ is in $U$.
(A3)~Let $x\in E_r$. Then $x$ was processed before $r$ (A2). At $x$'s turn, $a_x=Y_x$, every junk good, and $Y_r$ were available, and $\{b_x,c_x\}\subseteq J\cup\{Y_r\}$, so $R_x\subseteq G$ and $x$ is a leader by (I3). Hence each block contains at most one exposed agent; let $B^*$ be $r$'s block and $k^*$ its exposed agent, if any.
(A4)~For a frozen $x$, some $j\notin U$ has $Y_x\in N_j$, and by (B) $j$ is processed after $x$ in $x$'s block. Iterating while the current agent is frozen gives a need chain of distinct agents of $x$'s block ending at a terminal. Let $\tau(x)=x$ if $x\in T$, and otherwise the end of some need chain from $x$; it has $\capa\ge1$.
For $x\in E_r$, $\pi_x=\{b_x,c_x\}\cap J\ne\emptyset$, since at most one of $b_x,c_x$ is $Y_r$; for the same reason no exposed pair lies in $\base(r)$, the first condition of Lemma~\ref{lem:owner}.

The agents $\tau(x)$, $x\in E_r\setminus B^*$, lie in distinct blocks other than $B^*$, so they are distinct terminals other than $r$, and $S-\capa(r)\ge|E_r\setminus B^*|$. Let $H$ contain one good of each $\pi_x$, a common one when two meet. If $E_r\cap B^*=\emptyset$, then $|H|\le|E_r|\le S-\capa(r)$. Otherwise $E_r\cap B^*=\{k^*\}$, $k^*\ne r$ and $|E_r\setminus B^*|=|E_r|-1$. If $k^*\in T$, or some need chain from $k^*$ ends at a terminal other than $r$, then $\tau(k^*)$ can be chosen in $B^*\setminus\{r\}$, one more terminal, so $S-\capa(r)\ge|E_r|\ge|H|$. If two sets $\pi_x$ meet, then $|H|\le|E_r|-1\le S-\capa(r)$. If none of these holds, the bad case holds.\qed
\end{proof}

\begin{proof}[of Theorem~\ref{thm:B}]
Let $k=k^*$, $W=J\cup\{Y_r\}$; the new junk is $J'=W\setminus\{b_k,c_k\}$, and every pick other than $Y_r$ is still a pick in $P'$.
(a)~$P'$ is a pre-allocation: $Y'_{x_s}\in R_{x_s}$ since $x_s$ needed it, and $b_k,c_k\in W$ are neither picks of $P'$ nor goods $c_u$.
(b)~$\NA'\subseteq\NA$: $Y'_{x_s}\succ Y_{x_s}$ (or $Y_{x_s}=\none$), so $N'_{x_s}\subseteq N_{x_s}$; $k$ had $N_k=\emptyset$ and is now in $U'$.
(c)~Validity: $J'\subseteq W$, $J\cap\NA=\emptyset$ and $Y_r\notin\NA$ (A1), so $J'\cap\NA'=\emptyset$; for $u\in U$, $b_u,c_u\notin\NA\supseteq\NA'$; $b_k,c_k\in W$ are not in $\NA$.
(d)~A terminal $x\notin B^*$ keeps its pick and is a terminal of $P'$ with the same slots.
(e)~$E'_k\subseteq E_r\setminus\{k\}$, where $E'_k$ is taken in $P'$, with $J'\cup\{b_k,c_k\}\subseteq W$. An agent off the chain in $E'_k$ is in $E_r$. An interior chain agent $x_s$ ($1\le s<t$) is not: $a_{x_s}=Y'_{x_s}\succ Y_{x_s}$, so $Y_{x_s}\in\{b_{x_s},c_{x_s}\}$, and $Y_{x_s}\notin W$. Nor is $r$: $Y_r\ne a_r$. If $Y_r=\none$ or $Y_r=c_r$, then $b_r$ was picked by another agent (I1), so $b_r\notin W$. If $Y_r=b_r$, then (UT) applies to $r$ ($b_r\notin\NA$ by (A1)), so $c_r\notin J$, and $c_r\ne Y_r$, so $c_r\notin W$.
(f)~If some $x\in E'_k$ had $\{b_x,c_x\}=\{b_k,c_k\}$, then $x\in E_r\setminus\{k\}$ and $\pi_x=\pi_k\ne\emptyset$, against disjointness.
(g)~By (f), every $x\in E'_k$ has a good of $\{b_x,c_x\}$ in $J'$; let $H'$ contain one each. By (e), $|H'|\le|E_r|-1=|E_r\setminus B^*|$. The terminals $\tau(x)$, $x\in E_r\setminus B^*$, are distinct, outside $B^*$, and keep $\capa'\ge1$ by (d). So $|H'|\le S'=S'-\capa'(k)$, and Lemma~\ref{lem:owner} applies to $P'$ and $o=k$.\qed
\end{proof}

\paragraph{Size of the large bundle (written, not formalized).} In a valid pre-allocation each good of $\NA$ is the pick of exactly one frozen agent, so $|F|=|\NA|$; counting goods gives $\omega=m-2n+|\NA|$, and the owner's bundle has $\omega+2$ goods. The rotation does not increase $\NA$.

\section{Algorithm K3ALG}\label{app:algo}

Lists are in index order; ``first'' means smallest index. This is the program formalized as \lean{EFX.K3.algoC} (\texttt{proofs/k3\_algorithm.md} of~\cite{Repo}).

\begin{alltt}
Stage R.  A := all agents;  M := all goods.
loop:
  if |A| <= 1: the agent of A takes M; stop.
  if M is empty: stop.
  i := the first agent of A for which, with p the
       first good of M with the largest v_i(p),
       v_i(p) = 0 or v_i(M - p) <= v_i(p);
  if there is none: go to Stage L.
  remove i from A; if v_i(p) > 0, i takes p, and
       p is removed from M.
Stage L.  a_i, b_i, c_i := i's goods sorted by value.
  order: repeatedly the first unprocessed agent with
         at most two of its goods left, else the
         first unprocessed agent.
  Phase 1 in this order; upgrades (first eligible k).
  if |J| <= S: complete without owner; stop.
  r := the last agent of the order not in U;
  E := the agents exposed for r;
  H := one junk good of each exposed pair, with one
       common good for the first meeting pair;
  if |H| <= S - cap(r): complete with owner r,
       placing H first; stop.
  k := the agent of E in r's block;
  chain: scan forward from k, appending each agent
       outside U that needs the current pick, while
       the current agent is frozen;
  rotate; recompute J, cap and S;
  complete without owner if |J| <= S,
       else with owner k.
\end{alltt}

\noindent \emph{Cost table (Lean bound with $N=n+m+1$; finer bound written).} Stage R: $15N^2$ per round, $O(n^2m)$ in all. Rankings: $14N^2$. Processing order: $15N^3$. Phase~1's picks: $17N^3$. Upgrades: $36N^4$, $O(n^4+n^2m)$, the dominant term. Junk and slots: $O(N^3)$. Owner test: $20N^2+24N^3$. Blocks, $k^*$, chain, rotation, completion: $O(N^3)$. Output: $4N^2$. The total is at most $400N^4$ (\lean{EFX.K3.algoC_cost}), also at most $6400(n+m)^4$. On random instances the count is far below the bound (ratio $0.002$ at $n=30$).

\paragraph{Evidence.} Three implementations (a literal transcription of the Lean definitions, the fast Python version, and Lean's own evaluation of \lean{EFX.K3.algo}) agree on $200{,}000$ random instances (the first two) and on $675$ instances (all three), every output passes a raw \efxz{} check, and every Lean operation count is below the proved bound. The fast version takes $1.53$\,s and $1.58$\,s (medians) on random cores with $10^5$ agents and $1.5\times10^5$ and $2\times10^5$ goods.

\section{The Lean Formalization}\label{app:lean}

\paragraph{Trusted base.} The file \lean{lean/EFX/Model.lean}, copied from~\cite{MrdEfx} with the namespace renamed, defines:
\begin{alltt}
def finSum : (k : Nat) \(\to\) (Fin k \(\to\) Nat) \(\to\) Nat
  | 0, _ => 0
  | k+1, f => finSum k (fun i => f i.castSucc)
                + f (Fin.last k)
structure Inst where
  n : Nat
  m : Nat
  v : Fin n \(\to\) Fin m \(\to\) Nat
abbrev Alloc := Fin I.m \(\to\) Fin I.n
def bundleVal (X : I.Alloc) (i j : Fin I.n)
    (ex : Option (Fin I.m)) : Nat :=
  finSum I.m (fun g =>
    if X g = j \(\wedge\) ex \(\ne\) some g then I.v i g else 0)
def EFX0 (X : I.Alloc) : Prop :=
  \(\forall\) i j : Fin I.n, i \(\ne\) j \(\to\) \(\forall\) g : Fin I.m, X g = j \(\to\)
    I.bundleVal X i j (some g)
      \(\le\) I.bundleVal X i i none
def numRelevant (I : Inst) (i : Fin I.n) : Nat :=
  finSum I.m (fun g => if 0 < I.v i g then 1 else 0)
\end{alltt}

\paragraph{Main statements.}
\begin{alltt}
theorem EFX.target (I : Inst) (hn : 0 < I.n)
    (h : \(\forall\) i, numRelevant I i \(\le\) 3) :
    \(\exists\) X : I.Alloc, I.EFX0 X
theorem EFX.LB.corollaryD (I : Inst) (hn : 0 < I.n)
    (h3 : \(\forall\) i, numRelevant I i = 3)
    (hbal : \(\forall\) i g, 2 * I.v i g \(\le\) finSum I.m (I.v i)) :
    \(\exists\) X : I.Alloc, I.EFX0 X \(\wedge\) \(\exists\) o, \(\forall\) j, j \(\ne\) o \(\to\)
      finSum I.m (fun g => if X g = j then 1 else 0)
        \(\le\) 2
def EFX.K3.algo (I : Inst) (hn : 0 < I.n) : I.Alloc :=
  (algoC I hn).val
theorem EFX.K3.algo_efx0 (I : Inst) (hn : 0 < I.n)
    (h : \(\forall\) i, numRelevant I i \(\le\) 3) :
    I.EFX0 (algo I hn)
theorem EFX.K3.algoC_cost (I : Inst) (hn : 0 < I.n) :
    (algoC I hn).cost \(\le\) 400 * (I.n + I.m + 1) ^ 4
\end{alltt}
\texttt{EFX.target\_ordered} and \texttt{EFX.corollaryD\_ordered} state the same over any type with $0$, $+$ and $\le$ satisfying the axioms of a linearly ordered cancellative commutative monoid (class \texttt{EFX.OrderedValue}), with nonnegativity as a hypothesis, relevance as $v_i(g)>0$ and balance as $v_i(g)+v_i(g)\le v_i(M)$; at $\mathbb{N}$ they are exactly the statements above.

\paragraph{Theorems cited} (all in \texttt{lean/EFX/}; each has a \verb|#print axioms| certificate listing at most the three standard axioms; ledger rows of~\cite{Repo} in brackets):

\medskip
\noindent\begin{tabular}{@{}p{0.33\textwidth}p{0.63\textwidth}@{}}
\toprule
Result [ledger row] & Lean names (file in \texttt{lean/EFX/}) \\
\midrule
Thm.~\ref{thm:target} [T] & \lean{EFX.target} (\texttt{Target}); \lean{EFX.target_ordered} (\texttt{RealValues}) \\
Thm.~\ref{thm:D} [D] & \lean{EFX.LB.corollaryD} (\texttt{CorollaryD}); \lean{EFX.corollaryD_ordered} (\texttt{RealValues}) \\
Thm.~\ref{thm:algo} [K3.ALG] & \lean{EFX.K3.algo_efx0} (\texttt{K3Cost}) \\
Thm.~\ref{thm:algo} [K3.ALG.TIME] & \lean{EFX.K3.algoC_cost} (\texttt{K3CostBound}) \\
Lem.~\ref{lem:peel} [L2] & \lean{EFX.peel} (\texttt{Peeling}); \lean{EFX.peelEmpty}, \lean{EFX.peelR2} (\texttt{PeelingR2}) \\
Lem.~\ref{lem:junk} [L3] & \lean{EFX.Inst.exists_efx0_of_junk} (\texttt{Junk}) \\
Prop.~\ref{prop:core} [T] & \lean{EFX.core_reduction} (\texttt{Target}) \\
Lem.~\ref{lem:L12} [L12] & \lean{EFX.l12}, \lean{EFX.efx0_iff_of_agree} (\texttt{RealValues}) \\
Thm.~\ref{thm:sound} [S2.LB+] & \lean{EFX.LB.Valid.sound} (\texttt{PreAlloc}) \\
Thm.~\ref{thm:A} [S2.R] & \lean{EFX.LB.theoremA_invalid} (\texttt{OwnerR}) \\
Thm.~\ref{thm:B} [S2.LB+] & \lean{EFX.LB.theoremB} (\texttt{Rotation}) \\
Thm.~\ref{thm:C} [S2.LB+] & \lean{EFX.LB.lbPlusRun_sound}, \lean{EFX.LB.lbPlusOut_exists} (\texttt{LBPlus}) \\
Prop.~\ref{prop:O} [K3.OWNER] & \lean{EFX.LB.validOwner_iff} (\texttt{OwnerR}) \\
Audit [AUD] & \lean{Audit.target_audit}, \lean{Audit.corollaryD_audit} (\texttt{Audit}) \\
Second proof [K4.C4X.K3] & \lean{EFX.C4min.corollaryD_K3}, \lean{EFX.C4min.target_K3} (\texttt{K3Theorem}) \\
\bottomrule
\end{tabular}
\medskip

\noindent \lean{EFX.LB.lbPlusRun_sound} quantifies over all runs of \LBp{} (any Phase~1 order with R1 priority, any upgrade order, any need chain, any completion satisfying (OC)); \lean{EFX.LB.lbPlusOut_exists} shows that an output always exists. The audit's second kernel (nanoda) checked 4,186 declarations: the 726 named declarations of the library at the time (the natural-number development of Theorems~\ref{thm:target} and~\ref{thm:D}) and their dependencies. The files formalizing Theorems~\ref{thm:target}--\ref{thm:algo}, their proofs and the audit have about 9,700 lines.

\section{Computational Certification}\label{app:cert}

Each row is a ledger row of~\cite{Repo} with status CERTIFIED; certificates are stored allocations, re-checked without SAT from the raw \efxz{} definition, and the enumeration of connected cores is checked for completeness by a second, independent method.

\medskip
\noindent\begin{tabular}{@{}p{0.5\textwidth}p{0.38\textwidth}l@{}}
\toprule
Claim and scope & Size & Row \\
\midrule
\efxz{} exists, every $3$-relevant instance, $n\le6$ & $251+3{,}182$ connected cores & R1 \\
\efxz{} exists, every $3$-relevant instance, $n\le7$ & adds $n=7$ (rows R2, S2.N7) & R5 \\
Theorem~\ref{thm:D}, every core with $n\le6$, connected or not & certificates, components combined & R3 \\
Theorem~\ref{thm:D}, connected cores, $n=7$, $11\le m\le13$ & $37$, $541$, $3{,}103$ cores; $6^7$ profiles each & R2 \\
Theorem~\ref{thm:D}, connected cores, $n=8$, $m\in\{14,15\}$ & $1{,}232$ and $52$ cores & R4 \\
LB never fails, cores with $n\le6$ & $3{,}436$ connected, $131$ disconnected & S2.N6 \\
LB never fails, connected cores with $n=7$ & $41{,}170$ cores, $1.15\times10^{10}$ pairs & S2.N7 \\
LB never fails, connected cores, $n=8$, $m\ge13$ & $12{,}763$ cores, $2.14\times10^{10}$ pairs & S2.N8 \\
\bottomrule
\end{tabular}
\medskip

\noindent For the rows on LB, the certificates (LB's outputs) confirm Theorem~\ref{thm:D} in every labelling; that LB itself never fails holds for the labelling stored in the certificates, as LB breaks ties by index. Cores with $m=2n$ are covered by a two-own-goods argument (every agent holding two of its own goods is safe).

\paragraph{Evidence (not proof).} $5{,}844{,}263{,}904$ runs of a literal implementation of \LBp{} (every core with $n\le6$ under every ranking profile and every sequence of insertion choices; random choices on larger cores; random non-core instances), asserting every lemma of Section~\ref{sec:lbplus} and checking every output against the raw definition: $1{,}327{,}101$ rotations, $0$ failures. An independent Python implementation gives identical tallies for $n\le5$.

\paragraph{Refutations} (status REFUTED, with checkers):
\begin{itemize}
\item bundles of at most two goods suffice in a core: false for $n=3$, $m=5$ (Proposition~\ref{prop:neg}(a), $14$ of $216$ profiles); no core with $n=2$ fails;
\item bundles of at most two goods suffice when $m\le2n-2$: false, smallest $n=4$, $m=6$; at $n=6$, $m=10$, $57$ of the $211$ connected cores have a failing profile (confirmed by a second encoding);
\item the same for instances that are not cores: four identical agents with values $4,3,2$ on three goods, and three goods valued by nobody;
\item one bundle of three goods suffices: false for $n=4$, $m=7$ (Proposition~\ref{prop:neg}(b)); at $n=6$, $m=11$, $3$ of $25$ cores need four.
\end{itemize}

\section{Ordinality of Cores and the Counterexamples}\label{app:neg}

A good is \emph{alone} if it is the only good of its bundle.

\begin{lemma}[cases of safety]\label{lem:L5}
Let agent $i$ have $R_i=\{a,b,c\}$ with $v_i(a)>v_i(b)>v_i(c)>0$ and $v_i(a)<v_i(b)+v_i(c)$. Then $i$ satisfies the \efxz{} constraints iff one of: (T)~$i$ holds $a$, and either $i$ holds $b$ or $c$, or $b$ and $c$ are not together in a bundle of at least three goods; (P)~$i$ holds $b$ and $c$; (B)~$i$ holds $b$, and $a$ is alone; (C)~$i$ holds $c$, and $a$ and $b$ are alone; (E)~$a$, $b$ and $c$ are all alone.
\end{lemma}

\begin{proof}
The subset sums are ordered $0<c<b<a<b+c<a+c<a+b<a+b+c$ (writing $x$ for $v_i(x)$), so every constraint is decided by the ranking. Consider the goods of $R_i$ that $i$ holds. If it holds $a$ and one more, every other bundle meets $R_i$ in one good, worth $\le b<a+c$. If it holds $\{b,c\}$, the only threat is $a<b+c$. If it holds only $a$: bundles meeting $R_i$ in one good are harmless; $\{b,c\}$ as a bundle threatens $b<a$; $b,c$ together with another good threaten $b+c>a$. If it holds only $b$: $a$ not alone lies with $c$ (threat $a$) or with a good outside $R_i$ (threat $\ge a$), both $>b$; if $a$ is alone, $c$'s bundle threatens at most $c<b$. If it holds only $c$, it needs $a$ and $b$ alone by the same argument. If it holds none of $a,b,c$, any bundle with a good of $R_i$ and at least two goods threatens a positive amount, so all three must be alone.\qed
\end{proof}

\noindent The cases remain sufficient with ties, so tied cores reduce to strict ranking profiles (\texttt{proofs/lemmas.md}, L5).

\begin{proposition}[limits of the shape]\label{prop:neg}
Take strict rankings and strictly balanced values.
(a) Let agents $0,1,2$ rank $g_0\succ g_1\succ p_i$, where $p_i$ is private to agent $i$. No \efxz{} allocation has all bundles of at most two goods, while $\{g_0\},\{g_1\},\{p_0,p_1,p_2\}$ is \efxz{}.
(b) Let agents $0,1,2,3$ rank $g_0\succ g_2\succ p_0$, $g_0\succ g_2\succ p_1$, $g_1\succ g_2\succ p_2$, $g_1\succ g_2\succ p_3$. Every \efxz{} allocation with at most one bundle of more than two goods has bundle sizes $4,1,1,1$.
\end{proposition}

\noindent Exhaustive computation shows that no core with two agents needs a bundle of more than two goods, so (a) is a smallest example (ledger row X1 of~\cite{Repo}; \texttt{proofs/counterexamples.md} and \texttt{proofs/construction.md}, Section~1).

\begin{proof}[of Proposition~\ref{prop:neg}]
(a)~With five goods in three bundles of at most two goods, the sizes are $2,2,1$, so exactly one good is alone. Good $g_0$ is everyone's top and has one holder, so two agents do not hold their top. For them (T) fails, (C) and (E) need two or three alone goods, so each is in case (P) or (B), and both hold $g_1$: impossible. In $\{g_0\},\{g_1\},\{p_0,p_1,p_2\}$, the holder of $g_0$ is in case (T), the holder of $g_1$ in case (B), and the third agent in case (C).
(b)~Each of the tops $g_0,g_1$ has a claimant that does not hold it (a \emph{loser}). A loser is safe only by (P) or (B), which need it to hold $g_2$, or by (C) or (E), which need $g_2$ alone. The two losers are distinct and cannot both hold $g_2$, so $g_2$ is alone, and no agent is in case (P). A loser holding $\{g_2\}$ is in case (B) or (E) and needs its top alone; a loser not holding $g_2$ is in case (C) or (E) and needs its top and $g_2$ alone. So $g_0,g_1,g_2$ are alone, held by three agents, and the fourth bundle is $\{p_0,p_1,p_2,p_3\}$. The allocation giving $g_0$ to agent~0, $g_1$ to agent~2, $g_2$ to agent~1 and the private goods to agent~3 is \efxz{} (cases T, B, T, C).\qed
\end{proof}

\end{document}
